% TOP_PROBLEM: {"problemId": "problem.top500-omission-w045049-deligne-drinfeld-grt-one-freeness", "canonicalTitle": "Deligne–Drinfeld Freeness Conjecture for grt_1", "releaseRank": 448, "primaryDomain": "number_theory_arithmetic_geometry", "primaryDomainLabel": "Number theory & arithmetic geometry", "displayStatus": "open_with_solved_subcases", "releaseStatus": "open_with_solved_subcases"}
% !TeX program = lualatex
% Definition and sources only; this file is not an LLM solution attempt.
\documentclass[11pt]{article}
\usepackage[margin=25mm]{geometry}
\usepackage{fontspec}
\setmainfont{Latin Modern Roman}
\usepackage{amsmath,amssymb,mathtools}
\usepackage{unicode-math}
\setmathfont{Latin Modern Math}
\usepackage{xurl,hyperref}
\hypersetup{colorlinks=true,urlcolor=blue}
\setcounter{secnumdepth}{0}
\setlength{\parindent}{0pt}
\setlength{\parskip}{5pt}
\setlength{\emergencystretch}{3em}
\begin{document}
\section*{\texorpdfstring{Deligne–Drinfeld Freeness Conjecture for grt\_1}{Deligne–Drinfeld Freeness Conjecture for grt\_1}}\label{448}
\subsection{Definitions and mathematical statement}
Let $\mathfrak{grt}_1$ denote the weight-graded Grothendieck--Teichm\"uller Lie algebra over ${\mathbb{Q}}$. It is defined by Lie series in two generators satisfying Drinfeld's infinitesimal antisymmetry, three-cycle, and pentagon identities, with its Ihara bracket; the defining identities and completion conventions are those of the references in the cited paper. Weight is total degree in the two generators. This is not the full Grothendieck--Teichm\"uller algebra with its additional degree-zero grading direction.

The Deligne--Drinfeld freeness assertion is
\[
 \mathfrak{grt}_1\cong
 \mathbb L_{{\mathbb{Q}}}(\sigma_3,\sigma_5,\sigma_7,\ldots),
 \qquad \operatorname{wt}(\sigma_{2j+1})=2j+1\quad(j\ge1),
\]
where $\mathbb L$ denotes the free graded Lie algebra: all relations follow from bilinearity, antisymmetry, and Jacobi. For completed Lie algebras, both sides are completed by weight. The source also discusses larger conjectural chains of isomorphisms; those extra isomorphisms are not added to this entry's freeness target.
\par\smallskip\textit{Formulation and concept references:} \href{https://arxiv.org/html/2508.08081}{[S1]}.

\subsection{Short English statement}
Is the graded Grothendieck--Teichm\"uller Lie algebra freely generated by exactly one generator in each odd weight starting at three?
\subsection{Sources}
\begin{itemize}
\item[S1]Florian Naef and Thomas Willwacher, Numerical computation of linearized KV and the Deligne-Drinfeld and Broadhurst-Kreimer conjectures (2025 preprint).
\url{https://arxiv.org/html/2508.08081}.
\textit{Formulation source. Consulted 24 September 2026: Graded freeness question and distinction from the larger conjectural chain.}
\end{itemize}
\end{document}
